\section{Población y muestra}

\subsection{Población}
La población técnica comprende las hojas, cuadernos y fichas elaboradas en
Braille Grado 1 en español mediante regleta y punzón por estudiantes con
discapacidad visual de Moquegua e Ilo. La población pedagógica comprende a los
docentes de aula regular y especialistas de apoyo de ambas UGEL.

\subsection{Muestra}
Se conforma un corpus de 40 páginas manuscritas sobre cartulina escolar estándar
(140 a 160 g/m\textsuperscript{2}): 35 páginas controladas y 5 páginas con
producciones espontáneas. Cada página se captura bajo 12 combinaciones
fotométricas, lo que produce 480 fotografías originales. La partición se realiza
a nivel de página antes del aumento de datos para evitar fuga de información.

Para el componente pedagógico se prevé un máximo de 30 docentes de las UGEL Ilo y
Mariscal Nieto, justificado con un análisis de potencia a priori para una prueba
$t$ de muestras relacionadas ($\alpha = .05$, potencia $= .80$, $d = 0.55$).

\subsection{Criterios de inclusión y exclusión}
\begin{itemize}
  \item \textbf{Inclusión:} textos en Braille Grado 1 en español, elaborados con
  regleta y punzón, y docentes con estudiantes con ceguera en aula regular.
  \item \textbf{Exclusión:} Braille Grado 2, notación Nemeth, textos impresos a
  máquina y docentes sin vinculación con estudiantes con ceguera.
\end{itemize}
